\section{Background}
\label{sec:background}

\subsection{Threat model}
The attacker controls a share of the training data and plants a trigger with a target label. The defender receives the trained model and holds a small clean validation set drawn from the same distribution, here \HeldoutSize{} images, but never the training set and never a poisoned example. At deployment the defender sees one input at a time and must decide whether it carries a trigger before returning a prediction. This is the setting of the input-level detectors of \cref{sec:related}.

\subsection{The prediction shift phenomenon in backdoored networks}
PSBD \cite{li2025psbd} runs the model $k$ times on an input $x$ with dropout at rate $p$ active and compares each dropout prediction with the prediction without dropout. The paper defines the prediction shift indicator and the shift ratio of a set $\mathcal{D}$ as
\begin{equation}
  \begin{split}
    \ps(x) &= \mathbb{I}\big[\mathcal{Y}(x;\theta) \neq \mathcal{Y}(x;\theta')\big], \\
    \sigma(\mathcal{D}) &= \frac{1}{k|\mathcal{D}|}\sum_{x\in\mathcal{D}} \ps(x),
  \end{split}
\end{equation}
where $\theta$ are the weights, $\theta'$ the weights with dropout active and $\mathcal{Y}$ the predicted class. The detection statistic, called the prediction shift uncertainty, is the drop in the probability of the no-dropout class averaged over the $k$ passes,
\begin{equation}
  \psu(x) = \frac{1}{k}\sum_{i=1}^{k}
    \big[ P(\hat y \mid x;\theta) - P(\hat y \mid x;\theta'_i) \big],
\end{equation}
with $\hat y$ the class predicted without dropout. A poisoned input keeps its probability under dropout and has a low $\psu$.

The dropout rate is chosen so that the shift ratio of the clean validation set reaches \AdaptiveShiftTarget{}. The threshold is the \HeadlineQuantilePercent{} percentile of the validation set's $\psu$ at $k = 3$ passes. In a ResNet the paper applies dropout after the residual addition of each basic block and before its activation, and it explains the phenomenon by neuron bias, which \cref{sec:mechanism:phenomenon} tests on ViT.

\subsection{Perturbation sites in a ViT block}
A ViT-B/16 \cite{dosovitskiy2021an} splits a \ModelInputSize{} by \ModelInputSize{} image into \VitPatches{} patches of 16 by 16 pixels, embeds each as a \VitWidth{}-dimensional token, prepends a class token and runs \VitBlocks{} identical blocks. Each block is pre-norm \cite{xiong2020layernorm}. It applies a LayerNorm, multi-head self-attention over the \VitTokens{} tokens and a residual addition, then a second LayerNorm, a two-layer MLP applied to every token independently and a second residual addition. Neither addition is followed by an activation. The classifier reads the class token after a final LayerNorm. \Cref{fig:vit-block} shows one block with the sites where we inject a perturbation.

Attention is the only operation that mixes information across tokens. A perturbation before the attention LayerNorm therefore acts on each token before any mixing, a perturbation before the MLP LayerNorm acts after the block's mixing has happened, and a perturbation on the residual stream acts on the sum of everything the network has written so far.

\begin{figure}[t]
\centering
% One pre-norm ViT block with the perturbation sites. Every box has the same
% width so the residual skip, drawn left of the widest box, crosses none of them.
% A and D sit on the branch after the skip splits off, B on the branch output
% before the addition and C on the stream after each addition, which is where
% PSBD-RD applies dropout.
\begin{tikzpicture}[
  box/.style={draw, rounded corners=1pt, minimum width=30mm, text width=28mm, minimum height=5.5mm, font=\scriptsize, align=center},
  norm/.style={box, fill=black!6},
  op/.style={box, fill=blue!7},
  plus/.style={draw, circle, inner sep=0.8pt, font=\scriptsize, fill=white},
  site/.style={circle, fill=black!70, text=white, inner sep=0pt, minimum size=3.4mm, font=\bfseries\tiny},
  leader/.style={black!45, thin},
  flow/.style={-{Latex[length=1.4mm]}, thick}
]
\def\gap{5mm}
\def\skipx{-20mm}
\def\sitex{21mm}
\node[font=\scriptsize] (in) at (0,0) {residual stream, \VitTokens{} tokens $\times$ \VitWidth{}};
\coordinate[below=\gap of in] (split1);
\node[norm, below=\gap of split1] (ln1) {LayerNorm};
\node[op, below=\gap of ln1] (attn) {Multi-Head Attention};
\node[plus, below=\gap of attn] (add1) {$+$};
\coordinate[below=\gap of add1] (split2);
\node[norm, below=\gap of split2] (ln2) {LayerNorm};
\node[op, below=\gap of ln2] (mlp) {MLP, per token};
\node[plus, below=\gap of mlp] (add2) {$+$};
\node[font=\scriptsize, below=\gap of add2] (out) {to the next block};
\draw[flow] (in) -- (ln1);
\draw[flow] (ln1) -- (attn);
\draw[flow] (attn) -- (add1);
\draw[flow] (add1) -- (ln2);
\draw[flow] (ln2) -- (mlp);
\draw[flow] (mlp) -- (add2);
\draw[flow] (add2) -- (out);
\draw[flow] (split1) -- (\skipx,0 |- split1) |- (add1.west);
\draw[flow] (split2) -- (\skipx,0 |- split2) |- (add2.west);
\fill (split1) circle (0.5pt);
\fill (split2) circle (0.5pt);
% Each marker sits in 1 column right of the boxes, joined to the midpoint of
% the arrow it perturbs.
\coordinate (a) at ($(split1)!0.5!(ln1.north)$);
\coordinate (b) at ($(attn.south)!0.5!(add1.north)$);
\coordinate (c1) at ($(add1.south)!0.35!(split2)$);
\coordinate (d) at ($(split2)!0.5!(ln2.north)$);
\coordinate (c2) at ($(add2.south)!0.5!(out.north)$);
\foreach \point/\name in {a/A, b/B, c1/C, d/D, c2/C} {
  \draw[leader] (\point) -- (\sitex,0 |- \point);
  \node[site] at (\sitex,0 |- \point) {\name};
}
\end{tikzpicture}

\caption{One ViT block with the sites we perturb. A is the attention input before its LayerNorm, where PSBD-TM masks whole tokens. B is the attention branch output before its residual addition. C is the residual stream after each addition, where PSBD-RD applies dropout as our adaptation of the original PSBD site. D is the MLP input before its LayerNorm. The attention heads and the MLP hidden units are further sites we perturb.}
\label{fig:vit-block}
\end{figure}
